\subsection{NORMAL INTEGRAL SUBGROUPS}

Lemma 3. Let \(\mathbf{G}\) be a Lie group, \(\mathrm{H}_1\) and \(\mathrm{H}_2\) integral subgroups whose topology admits a countable base and \(g \in \mathbf{G}\) . Then

\[
g \mathrm{H} _ {1} g ^ {- 1} = \mathrm{H} _ {2} \Leftrightarrow (\operatorname{Ad} g) (\mathrm{L} (\mathrm{H} _ {1})) = \mathrm{L} (\mathrm{H} _ {2}).
\]

\(\operatorname{Ad} g = \mathrm{T}_{e}(\operatorname{Int} g)\) . Hence, by transport of structure, \((\operatorname{Int} g)(\mathrm{H}_{1})\) has Lie algebra \((\operatorname{Ad} g)(\mathrm{L}(\mathrm{H}_{1}))\) . As \(\mathrm{H}_{1}\) and \(\mathrm{H}_{2}\) have countable bases, to say that the sets \(\mathrm{H}_{2}\) and \((\operatorname{Int} g)(\mathrm{H}_{1})\) are equal amounts to saying that the integral subgroups \(\mathrm{H}_{2}\) and \((\operatorname{Int} g)(\mathrm{H}_{1})\) are equal (no. 2, Corollary 3 to Proposition 3). Then the lemma follows from Theorem 2 (i).

PROPOSITION 14. Let G be a Lie group and H an integral subgroup whose topology admits a countable base. The following conditions are equivalent:

\begin{enumerate}
\def\labelenumi{(\roman{enumi})}
\item
  H is normal in G;
\item
  L(H) is invariant under Ad(G).
\end{enumerate}

If further G is connected, these conditions are equivalent to the following:

\begin{enumerate}
\def\labelenumi{(\roman{enumi})}
\setcounter{enumi}{2}
\tightlist
\item
  L(H) is an ideal of L(G).
\end{enumerate}

If further \(\mathbf{G}\) is simply connected and \(\mathbf{L}(\mathbf{H})\) is of finite codimension in \(\mathbf{L}(\mathbf{G})\) , these conditions imply that \(\mathbf{H}\) is a Lie subgroup of \(\mathbf{G}\) and that \(\mathbf{G} / \mathbf{H}\) is simply connected.

The equivalence (i) \(\Leftrightarrow\) (ii) follows from Lemma 3. If further G is connected, condition (ii) is equivalent to saying that L(H) is stable under ad L(G) (no. 5, Corollary 1 to Proposition 13 and § 3, no. 12, Proposition 44).

Suppose that G is simply connected and that L(H) is an ideal of finite codimension in L(G). By Theorem 3 of no. 3, there exists a simply connected Lie group G' such that L(G') is isomorphic to L(G)/L(H). There exists a continuous morphism f of L(G) onto L(G') with kernel L(H). By Theorem 1 of no. 1, there exists a morphism \(\phi\) of G into G' such that L( \(\phi\) ) = f.~This morphism is a submersion and hence its kernel N is a Lie subgroup of G such that L(N) = Ker f = L(H). Hence H is the identity component of N and is therefore a Lie subgroup of G. Let \(\psi\) be the morphism of G/H into G' derived from \(\phi\) when passing to the quotient. This morphism is étale; since G is simply connected, \(\psi\) is an isomorphism of G/H onto G'.

COROLLARY 1. Let \(G\) be a finite-dimensional simply connected Lie group. Let \(m, h\) be Lie subalgebras of \(L(G)\) such that \(L(G)\) is the semi-direct product of \(m\) by \(h\) . Let \(M, H\) be the corresponding integral subalgebras of \(G\) . Then \(M\) and \(H\) are simply connected Lie subgroups of \(G\) and, as a Lie group, \(G\) is the semi-direct product of \(M\) by \(H\) .

By Proposition 14, H is a normal Lie subgroup of G and the Lie group G/H is simply connected. Let \(\pi\) be the canonical morphism of G onto G/H. There exists a morphism \(\theta\) of G/H into M such that L(θ) is the canonical isomorphism of L(G)/L(H) = L(G)/h onto L(M) = m. Then

\[
\mathrm{L} (\pi \circ \theta) = \mathrm{L} (\pi) \circ \mathrm{L} (\theta) = \mathrm{Id} _ {\mathrm{L} (\mathrm{G} / \mathrm{H})}
\]

and hence \(\pi \circ \theta = \mathrm{Id}_{\mathrm{G}/\mathrm{H}}\) . By no. 1, Corollary 1 to Proposition 1, \(\theta (\mathrm{G / H}) = \mathbf{M}\) . By Proposition 8 of § 1, no. 4, M is a Lie subgroup of G and the Lie group G is the semi-direct product of M by H.

COROLLARY 2. Let G be a finite-dimensional simply connected Lie group, H a normal connected Lie subgroup of G and \(\pi\) the canonical morphism of G onto G/H.

\begin{enumerate}
\def\labelenumi{(\roman{enumi})}
\item
  There exists an analytic mapping \(\rho\) of G/H into G such that \(\pi \circ \rho = \mathrm{Id}_{\mathrm{G / H}}\) .
\item
  For every mapping \(\rho\) with the properties of (i), the mapping \((h, m) \mapsto h\rho(m)\) of \(\mathbf{H} \times (\mathbf{G}/\mathbf{H})\) into \(\mathbf{G}\) is an isomorphism of analytic manifolds.
\item
  H and G/H are simply connected.
\end{enumerate}

Let \(n = \dim G - \dim H\) . The corollary is obvious for \(n = 0\) . We argue by induction on \(n\) .

Suppose that there exists an ideal of L(G) containing L(H) distinct from L(G) and L(H). Let H' be the corresponding connected Lie subgroup of G. Let \(\pi_{1}: G \to G/H'\) and \(\pi_{2}: H' \to H'/H\) be the canonical morphisms. By the induction hypothesis, there exist analytic mappings \(\rho_{1}: G/H' \to G\) , \(\rho_{2}: H'/H \to H'\) such that \(\pi_{1} \circ \rho_{1} = Id_{G/H'}\) , \(\pi_{2} \circ \rho_{2} = Id_{H'/H}\) . Let

\[
\pi_ {3} \colon \mathrm{G} / \mathrm{H} \rightarrow \mathrm{G} / \mathrm{H} ^ {\prime}
\]

be the canonical morphism. If \(x \in \mathrm{G} / \mathrm{H}\) and \(y\) is a representative of \(x\) in \(\mathrm{G}, y\) and \(\rho_1(\pi_3(x))\) have the same canonical image modulo \(\mathrm{H}'\) and hence \(x^{-1}\pi (\rho_1(\pi_3(x))) \in \mathrm{H}' / \mathrm{H}\) . Let

\[
\rho (x) = \rho_ {1} (\pi_ {3} (x)) \rho_ {2} (\pi (\rho_ {1} (\pi_ {3} (x))) ^ {- 1} x) \in G.
\]

Clearly \(\rho\) is an analytic mapping of G/H into G and

\[
\begin{array}{r l} \pi (\rho (x)) & = \pi (\rho_ {1} (\pi_ {3} (x))) \pi_ {2} (\rho_ {2} (\pi (\rho_ {1} (\pi_ {3} (x))) ^ {- 1} x)) \\ & = \pi (\rho_ {1} (\pi_ {3} (x))) \pi (\rho_ {1} (\pi_ {3} (x))) ^ {- 1} x = x. \end{array}
\]

If now the only ideals of \(\mathbf{L}(\mathbf{G})\) containing \(\mathbf{L}(\mathbf{H})\) are \(\mathbf{L}(\mathbf{G})\) and \(\mathbf{L}(\mathbf{H})\) , the Lie algebra \(\mathrm{L(G) / L(H)}\) is either 1-dimensional or simple. In both cases, \(\mathrm{L(G)}\) is the semi-direct product of a subalgebra by \(\mathrm{L(H)}\) ; this is obvious in the first case and in the second this follows from Chapter I, § 6, Corollary 3 to Theorem 5. Assertion (i) then follows from Corollary 1.

Assertion (ii) is obvious. Assertion (iii) follows from (i) and (ii).

The conclusions of Corollary 2 are no longer necessarily true when G is infinite-dimensional or when H is not normal (Exercises 8 and 15).

COROLLARY 3. Let \(\mathbf{G}\) be a finite-dimensional connected Lie group and \(\mathbf{H}\) a normal connected Lie subgroup of \(\mathbf{G}\) . The canonical morphism of \(\pi_1(\mathbf{H})\) into \(\pi_1(\mathbf{G})\) is injective.

Let \(G_{1}\) be the universal covering of G and \(\lambda\) the canonical mapping of \(G_{1}\) onto G. The Lie algebra of \(G_{1}\) is identified with L(G). The Lie subgroup \(\lambda^{-1}(H)\) of \(G_{1}\) is normal in \(G_{1}\) and its Lie algebra is L(H). Let \(H_{1}\) be the identity component of \(\lambda^{-1}(H)\) and \(\lambda_{1} = \lambda | H_{1}\) . Then \(\mathrm{L}(H_{1}) = \mathrm{L}(H)\) and hence \(\lambda_{1}\) is an étale morphism of \(H_{1}\) onto H. On the other hand, \(H_{1}\) is simply connected (Corollary 2) and hence is identified with the universal covering of H. Then, by General Topology, Chapter XI, the canonical morphism of \(\pi_{1}(H)\) into \(\pi_{1}(G)\) is identified with the canonical injection of Ker \(\lambda_{1}\) into Ker \(\lambda\) .

